\chapter{Cel i zakres pracy}

\section{Cel pracy}
Celem pracy jest zaprojektowanie i implementacja aplikacji mobilnej w środowisku Android Studio, zintegrowanej z bazą danych MySQL, umożliwiającej odczyt i prezentację danych temperatury pochodzących z czujnika DHT11 podłączonego do modułu ESP32 WROOM.

\section{Cele szczegółowe}
W~ramach realizacji projektu wyznaczono następujące cele szczegółowe:
\begin{itemize}
    \item konfiguracja modułu ESP32 WROOM do odczytu danych z czujnika DHT11,
    \item implementacja mechanizmu zapisu danych (data, godzina, temperatura) do bazy MySQL,
    \item utworzenie lokalnego serwera bazy danych przy użyciu XAMPP,
    \item opracowanie interfejsu API (np. w PHP) do komunikacji aplikacji z bazą danych,
    \item zaprojektowanie aplikacji mobilnej w Android Studio,
    \item implementacja mechanizmu pobierania danych z bazy (REST API, JSON),
    \item wyświetlanie aktualnej temperatury oraz dziennika pomiarów w aplikacji,
    \item przeprowadzenie testów funkcjonalnych i niefunkcjonalnych.
\end{itemize}

\section{Zakres pracy}
Zakres pracy obejmuje:
\begin{itemize}
    \item konfigurację sprzętową (ESP32 + DHT11),
    \item konfigurację środowiska XAMPP i MySQL,
    \item implementację RestAPI,
    \item stworzenie aplikacji mobilnej Android,
    \item testowanie systemu.
\end{itemize}
Projekt nie obejmuje wdrożenia systemu w środowisku produkcyjnym ani zabezpieczeń klasy enterprise.

\section{Struktura pracy}
Praca została podzielona na pięć rozdziałów. Rozdział pierwszy przedstawia cel oraz zakres.
Rozdział drugi opisuje problematykę projektu i uzasadnia wybór tematu. Rozdział trzeci
zawiera analizę wymagań oraz projekt systemu. Rozdział czwarty prezentuje implementację
i~sposób uruchamiania rozwiązania. Rozdział piąty obejmuje testy i ocenę rezultatów.
